% =====================================================================================
% Chapter 1 · Potential outcomes, confounding, and DAGs
%
% Built from notebooks/00_foundations.ipynb (§1–§6, §8, §9). NOT ONE NUMBER IN THIS FILE
% IS TYPED BY HAND: every figure in the prose is a macro from build/macros.tex, generated
% by cmp.report from the executed notebook (CMP_FAST=0, CMP_REAL=1). The only literals
% below are (i) the constants of the data-generating models in §1.2, which are *definitions*
% rather than results, (ii) the conventional thresholds (the 0.1 SMD bar, the [0.05, 0.95]
% propensity clip, the 90 % interval level), which are *conventions* stated in the text, and
% (iii) the published LaLonde experimental benchmark, which is emitted as a macro and
% flagged, in the notebook and here, as cited rather than estimated.
% =====================================================================================
\chapter{Potential outcomes, confounding, and DAGs}
\label{ch:foundations}

\begin{quote}\small
\emph{A retailer's dashboard reports that emailed customers spend \eur{\nbZeroObsNaive} more
than un-emailed ones, against a contact cost of \eur{\nbZeroCost}. The email therefore pays,
and the recommendation writes itself. It is wrong. The true average effect in this world is
\eur{\nbZeroTrueAte}, and the gap is not noise: the confounded estimate sits
\nbZeroObsSeDist{} standard errors from the truth, while the randomized one sits
\nbZeroRandSeDist{}. Exactly \eur{\nbZeroSelBias} of the gap is a statement about \emph{who}
was emailed rather than about what the email did, and acting on the dashboard burns
\eur{\nbZeroPolBurn} on a base of \nbZeroNCustomers{} customers. This chapter builds the
vocabulary that makes that sentence precise --- potential outcomes, confounding, backdoor
paths, colliders --- and then climbs the estimator ladder that repairs the number, from a
naive difference of means through regression adjustment and inverse-propensity weighting to
the doubly-robust estimator, each with the standard error its data structure demands. It
closes on LaLonde's job-training data, where the naive comparison says training destroyed
\usd{\nbZeroLlNaive} of earnings and adjustment moves it \usd{\nbZeroLlSwing} --- across zero
--- to within \usd{\nbZeroLlShortfall} of the randomized benchmark.}
\end{quote}

% =====================================================================================
\section{The decision}
\label{sec:fnd:decision}

A retailer has a customer file and an email programme. Each contact costs \eur{\nbZeroCost}
to print, post and process. The analytics team pulls the obvious number --- average spend
among customers who were emailed, minus average spend among those who were not --- and it
comes to \eur{\nbZeroObsNaive}. The email clears its cost by a comfortable margin. Scale it.

That recommendation is the most expensive sentence in marketing analytics, and it is wrong
for a reason that has nothing to do with the arithmetic. The arithmetic is correct: the
emailed customers really did spend \eur{\nbZeroObsNaive} more. What is wrong is the
\emph{claim} smuggled into the comparison --- that the un-emailed customers tell us what the
emailed ones \emph{would have spent} had no email been sent. They do not, because nobody
flipped a coin. A marketer chose whom to email, and chose the loyal, engaged, recently-active
customers, who spend more anyway.

The chapter's central quantity is therefore not a number but a \emph{missing} number: what
the treated would have done untreated. Everything that follows --- the notation, the graphs,
the four estimators, the diagnostics --- exists to reconstruct it, and to state honestly the
assumptions under which the reconstruction is licensed. \Cref{sec:fnd:euro} prices what
getting it wrong costs, and the figure it arrives at, \eur{\nbZeroPolBurn} burned per
\nbZeroNCustomers{} contacts, is the reason this is not a philosophy chapter.

% =====================================================================================
\section{The data-generating model}
\label{sec:fnd:dgm}

A causal method cannot be graded on real data, because on real data the counterfactual is
missing --- that is the entire problem. So the methods are first graded against worlds whose
answers are known by construction, and only then trusted on a world whose answer is not.
\Cref{sec:fnd:real} performs the second half of that bargain on LaLonde's job-training data;
this section builds the first.

\subsection*{The customer file}

Each of the \nbZeroNCustomers{} customers carries features $x$ --- recency $R$, frequency
$F$, monetary value $M$, tenure and engagement $E$ --- a baseline spend $\mu_0(x)$, and a
\emph{planted per-customer effect} $\tau(x)$. Observed spend is
%
\begin{equation}
  Y \;=\; \mu_0(x) \;+\; \tau(x)\,T \;+\; \varepsilon,
  \qquad \varepsilon \sim \mathcal N\!\left(0,\ 8^{2}\right),
  \label{eq:fnd:uplift}
\end{equation}
%
with $T \in \{0,1\}$ the email. The simulator is run twice, and \emph{only the assignment rule
differs}:
%
\begin{align}
  \text{randomized:}\quad     & T \sim \text{Bernoulli}(\tfrac12), \label{eq:fnd:rct}\\
  \text{observational:}\quad  & T \sim \text{Bernoulli}\Big(\sigma\big(1.8\,(E - 0.4)
     + 1.2\,\tfrac{150 - R}{150} + 0.15\,(F - 5)\big)\Big), \label{eq:fnd:obs}
\end{align}
%
where $\sigma(z) = 1/(1 + e^{-z})$ is the logistic function, which turns any score into a
probability. \Cref{eq:fnd:obs} is a caricature of a real marketing department: engaged,
recently-active, frequent buyers are the ones who get emailed --- and, through $\mu_0(x)$,
they are the ones who would have spent more regardless. The planted average effect is
\eur{\nbZeroTrueAte} in both worlds. Nothing about the \emph{effect} changes between
\cref{eq:fnd:rct} and \cref{eq:fnd:obs}; only the mechanism that assigns it changes, and that
alone is enough to move the naive estimate from \eur{\nbZeroRandNaive} to
\eur{\nbZeroObsNaive} (\cref{fig:fnd:confounding}).

\input{build/floats/nb00_confounding}

\subsection*{The one-confounder world}

The estimator ladder of \cref{sec:fnd:classical} needs a world simple enough that the reader
can hold the whole causal structure in mind, and it gets one. A single observed confounder,
loyalty $L \sim \mathcal N(0,1)$, drives both the send and the spend; an independent
unobserved trait $G \sim \mathcal N(0,1)$ adds noise to spend but never touches assignment,
and is therefore \emph{not} a confounder:
%
\begin{align}
  T &\;\sim\; \text{Bernoulli}\big(\sigma(0.8\,L)\big), \label{eq:fnd:dag_t}\\
  Y &\;=\; 20 \;+\; 5\,L \;+\; 6\,T \;+\; 3\,G \;+\; \varepsilon,
       \qquad \varepsilon \sim \mathcal N\!\left(0,\ 4^{2}\right). \label{eq:fnd:dag_y}
\end{align}
%
The planted effect is exactly \eur{\nbZeroTrueDag}, and it is the grade every estimator in
this chapter --- and the posterior of Chapter~2 --- is marked against. The coefficient
$5$ on $L$ in \cref{eq:fnd:dag_y} and the coefficient $0.8$ on $L$ in \cref{eq:fnd:dag_t} are
the two halves of the confounding: loyalty raises the chance of being emailed \emph{and}
raises spend. Remove either and the naive estimator would be fine.

\subsection*{The zero-effect world, for the collider}

\Cref{sec:fnd:collider} needs a world in which the true effect is \emph{exactly zero} and the
treatment is randomized, so that any effect an estimator reports is manufactured rather than
merely mismeasured:
%
\begin{equation}
  T \sim \text{Bernoulli}(\tfrac12), \qquad
  Y = 20 + 0 \cdot T + \varepsilon, \quad \varepsilon \sim \mathcal N(0, 5^2), \qquad
  C = \mathbf 1\big\{\,1.2\,T + 1.2\,\tilde Y + \eta > 0.8\,\big\},
  \label{eq:fnd:collider}
\end{equation}
%
with $\eta \sim \mathcal N(0,1)$ and $\tilde Y$ spend standardized to unit variance. Here $C$
is ``responded to the campaign'' --- more likely if the customer was emailed \emph{and} more
likely if the customer spends a lot. It is a common effect of the treatment and the outcome,
and \cref{eq:fnd:collider} is the smallest machine that manufactures a causal effect out of
nothing.

% =====================================================================================
\section{Identification}
\label{sec:fnd:ident}

\subsection*{Potential outcomes, and the fundamental problem}

Fix a unit $i$ --- a customer --- and a binary treatment $T_i$. In the notation of
\citet{neyman1923} and \citet{rubin1974}, the unit carries \emph{two} potential outcomes:
$Y_i(1)$, its spend if emailed, and $Y_i(0)$, its spend if not. Both exist before anything is
done; the treatment decides which one is revealed. What is observed is
%
\begin{equation}
  Y_i \;=\; T_i\,Y_i(1) \;+\; (1 - T_i)\,Y_i(0),
  \label{eq:fnd:consistency}
\end{equation}
%
which is called \emph{consistency}: the observed outcome equals the potential outcome
corresponding to the treatment actually received. The individual causal effect is the contrast
$Y_i(1) - Y_i(0)$, and it is \textbf{never observable for any unit}, because one of its two
terms is always missing. This is the \emph{fundamental problem of causal inference}
\citep{holland1986}, and its most useful reading is the one \citet{ding2018} make explicit:
causal inference is a \emph{missing-data} problem. Half of every unit's table was never filled
in, and every method in this book is a principled way of imputing it.

Because individual effects are unavailable, decisions run on averages, and \emph{which} average
is a choice that must be made before a method is chosen. Four appear in this book:
%
\begin{align}
  \text{ATE} &\;=\; \E\big[Y(1) - Y(0)\big]
     &&\text{the effect on a randomly chosen customer,} \label{eq:fnd:ate}\\
  \text{ATT} &\;=\; \E\big[Y(1) - Y(0) \,\big|\, T = 1\big]
     &&\text{the effect on those actually treated,} \label{eq:fnd:att}\\
  \tau(x) &\;=\; \E\big[Y(1) - Y(0) \,\big|\, X = x\big]
     &&\text{the CATE: the effect for customers like } x, \label{eq:fnd:cate}\\
  \text{LATE} &\;=\; \E\big[Y(1) - Y(0) \,\big|\, \text{compliers}\big]
     &&\text{the effect on the units an instrument moves.} \label{eq:fnd:late}
\end{align}
%
These answer different commercial questions about the same campaign. The ATE prices \emph{roll
it out to everyone?}; the ATT audits \emph{what did our past sends earn?}; the CATE decides
\emph{who next?}; the LATE is the honest scope when exposure could only be encouraged, not
assigned. They are not interchangeable, and \cref{fig:fnd:tau} shows why the distinction is
commercial rather than terminological: the ATE of \eur{\nbZeroTrueAte} sits below the
\eur{\nbZeroCost} cost, so an average-level rule says \emph{email no one} --- while
\pc{\nbZeroShareAboveCost} of customers individually clear the cost bar and
\pc{\nbZeroShareNegative} are actively \emph{harmed} by the email. One number, opposite
per-customer instructions. Recovering $\tau(x)$ is the subject of Chapter~3; naming the
estimand is the subject of this paragraph.

\input{build/floats/nb00_tau_spread}

\subsection*{The naive difference, decomposed}

Now the central algebraic fact of the chapter. Take the number the dashboard computes, apply
consistency \cref{eq:fnd:consistency}, and add and subtract the one quantity nobody
observes --- the mean untreated outcome \emph{among the treated}, $\E[Y(0) \mid T=1]$:
%
\begin{equation}
  \underbrace{\E[Y \mid T{=}1] - \E[Y \mid T{=}0]}_{\text{the naive difference (observable)}}
  \;=\;
  \underbrace{\E[Y(1) - Y(0) \mid T{=}1]}_{\text{ATT}}
  \;+\;
  \underbrace{\E[Y(0) \mid T{=}1] - \E[Y(0) \mid T{=}0]}_{\text{selection bias}} .
  \label{eq:fnd:decomp}
\end{equation}
%
\Cref{eq:fnd:decomp} is an identity, not an approximation, and it is the whole subject in one
line. The naive difference is a causal effect \emph{plus} a statement about who was chosen.
The simulator can verify it because it knows every customer's unobserved branch
$Y(0) = Y - \tau T$: the naive \eur{\nbZeroNaiveObs} decomposes exactly into an ATT of
\eur{\nbZeroAtt} and a selection bias of \eur{\nbZeroSelBias}. And the ATT itself exceeds the
ATE of \eur{\nbZeroAte} by \eur{\nbZeroTargetingGap}, because past targeting favoured the very
customers on whom the email works best. So the \eur{\nbZeroObsNaive} the dashboard reports and
the \eur{\nbZeroTrueAte} the business needs are separated by two distinct errors, and only one
of them is bias in the usual sense.

\subsection*{Randomization, and what it buys}

Under \cref{eq:fnd:rct} the coin flip makes $T$ independent of \emph{everything} a customer
carries into the experiment, measured or not --- in particular of $Y(0)$. The selection term in
\cref{eq:fnd:decomp} is then \textbf{exactly zero}, the ATT equals the ATE, and the naive
difference of means becomes an unbiased estimate of the causal effect. That is the entire proof
behind the green bar of \cref{fig:fnd:confounding}, and it is why an A/B test is worth more than
any estimator in this book.

The residual gap in the randomized arm --- \eur{\nbZeroRandNaive} against a truth of
\eur{\nbZeroTrueAte} --- is \emph{noise}: \nbZeroRandSeDist{} standard errors, and it shrinks
like $1/\sqrt{n}$. The gap in the observational arm --- \eur{\nbZeroObsNaive} --- is
\emph{bias}: \nbZeroObsSeDist{} standard errors, and it does not shrink at all. More data makes
a confounded estimate more precisely wrong. This distinction is the reason the chapter's central
diagnostic (\cref{fig:fnd:ladder}, right panel) is a picture of \emph{repeated samples} rather
than of one.

\subsection*{The assumptions, when there is no coin flip}

Without randomization the effect is not identified by the data alone. It is identified by the
data \emph{plus assumptions}, and those assumptions are the price of admission. Four are stated
here; each is either tested later in the chapter or explicitly declared untestable --- an
assumption with no accompanying test is a wish.

\begin{assumption}[SUTVA: no interference, one version of treatment]\label{as:fnd:sutva}
Unit $i$'s potential outcomes depend on $i$'s own treatment only, not on anyone else's, and
``emailed'' means one well-defined thing \citep{rubin1980, imbens2015}. In a customer file this
fails if emailed customers talk to un-emailed ones, or if ``the email'' was in fact three
creatives. \textbf{Untestable from the data}, and named as such: it is discharged by knowing how
the campaign ran, not by a statistic. Chapter~9 plants a violation in a simulator, which is
the only place one can be seen.
\end{assumption}

\begin{assumption}[Unconfoundedness]\label{as:fnd:unconf}
Conditional on the measured features, treatment is as good as randomly assigned:
%
\begin{equation}
  \{\,Y(0),\, Y(1)\,\} \;\perp\; T \;\;\big|\;\; X .
  \label{eq:fnd:unconf}
\end{equation}
%
Within customers alike on $X$, who got the email carries no further information about what they
would have spent. This is the load-bearing assumption of the entire chapter, and it is
\textbf{untestable}: it is a statement about the potential outcomes that were never observed.
\Cref{sec:fnd:wrong} does the only honest thing available --- it plants a violation and measures
what it costs (\cref{fig:fnd:hidden}) --- and Chapter~3 prices the strength a hidden
confounder would need to overturn a decision.
\end{assumption}

\begin{assumption}[Positivity, or overlap]\label{as:fnd:overlap}
Every kind of customer could have gone either way:
%
\begin{equation}
  0 \;<\; e(x) \;<\; 1, \qquad e(x) \;=\; \Prob(T = 1 \mid X = x),
  \label{eq:fnd:overlap}
\end{equation}
%
where $e(x)$ is the \emph{propensity score} \citep{rosenbaum1983}. If some region of feature
space is always emailed, there is no like-for-like comparison to be made there and no
reweighting can invent one. Unlike \cref{as:fnd:unconf}, this one is \textbf{testable}, because
$e(x)$ is estimable from observed data. It is tested in \cref{sec:fnd:valid}, where it passes
(\cref{fig:fnd:overlap}), and again in \cref{sec:fnd:real}, where it fails badly
(\cref{fig:fnd:lalonde}).
\end{assumption}

\begin{assumption}[A valid adjustment set]\label{as:fnd:backdoor}
The variables conditioned on block every backdoor path from treatment to outcome, and open no
new ones: $X$ satisfies Pearl's \emph{backdoor criterion} \citep{pearl1995, pearl2009}. In
practice this means $X$ contains the common causes and \emph{no descendants of the treatment} ---
no mediators, no colliders. It is \textbf{partially testable}: the graph itself is an
assumption, but the consequence of getting it wrong is measurable, and \cref{sec:fnd:collider}
measures it on a treatment whose true effect is exactly zero.
\end{assumption}

Under \crefrange{as:fnd:sutva}{as:fnd:backdoor}, the ATE is identified by the
\emph{adjustment formula}, also called the g-formula:
%
\begin{equation}
  \text{ATE} \;=\; \E_X\Big[\ \E\big[Y \mid T{=}1,\, X\big] \;-\; \E\big[Y \mid T{=}0,\, X\big]\ \Big].
  \label{eq:fnd:gformula}
\end{equation}
%
Read \cref{eq:fnd:gformula} as an instruction rather than a formula: \emph{compare like with
like within strata of $X$, then average the within-stratum contrasts over the population
distribution of $X$} --- crucially, over the \emph{population} mix, not over the mix inside one
arm, which is exactly the substitution the naive difference makes.

\Cref{tab:fnd:strata} runs \cref{eq:fnd:gformula} by hand, once, before any model is fitted.
Cutting loyalty into five quintiles and computing the emailed-minus-not difference \emph{within}
each gives roughly the planted \eur{\nbZeroTrueDag} in every bin; the strata-weighted average is
\eur{\nbZeroStratAte}. The pooled contrast on the same data is \eur{\nbZeroPooledNaive}. The
\emph{\% emailed} column is the whole explanation: the emailed share climbs from
\pc{\nbZeroStratLoEmailed} in the least loyal quintile to \pc{\nbZeroStratHiEmailed} in the most
loyal one, so the pooled number is comparing loyal emailed customers with casual controls. Every
rung of the ladder in \cref{sec:fnd:classical} is a smarter way of doing exactly this
stratification.

\input{build/tables/nb00_strata}

% =====================================================================================
\section{The graph: seeing, doing, and what to condition on}
\label{sec:fnd:collider}

\subsection*{Seeing is not doing}

\emph{Observing} $X = x$ selects the units that happened to arrive at $x$, carrying every reason
they got there. \emph{Doing} $do(X = x)$ reaches in and sets it \citep{pearl2009}. Causal
questions are about doing. Confounding is precisely the disagreement between the two: in
\cref{tab:fnd:strata}, seeing gives \eur{\nbZeroPooledNaive} and doing is worth
\eur{\nbZeroTrueDag}. When the backdoor variables are measured, the backdoor adjustment formula
converts one into the other,
%
\begin{equation}
  \Prob\big(y \mid do(T{=}t)\big) \;=\; \sum_x \Prob\big(y \mid T{=}t,\, X{=}x\big)\, \Prob(X{=}x),
  \label{eq:fnd:backdoor}
\end{equation}
%
which is \cref{eq:fnd:gformula} written for distributions rather than means. \Cref{tab:fnd:strata}
\emph{is} \cref{eq:fnd:backdoor}, executed by hand on five strata.

\subsection*{Three graphs, and only one calls for adjustment}

A directed acyclic graph draws each variable as a node and each direct cause as an arrow. It is
not decoration: it is the object that decides what to control for, and it must be committed to
before the data are touched \citep{greenland1999}. \Cref{fig:fnd:dags} shows the three
configurations that matter, and the instinct to ``control for everything we have'' gets two of
the three wrong.

\input{build/floats/nb00_dags}

\begin{description}[leftmargin=*, style=nextline]
  \item[Confounder --- a backdoor path: $T \leftarrow L \rightarrow Y$.]
    Association flows from $T$ to $Y$ through a path with an arrow \emph{into} $T$. This is
    non-causal association, and conditioning on $L$ blocks it. Adjust.
  \item[Mediator --- on the causal path: $T \rightarrow M \rightarrow Y$.]
    Whether the email was \emph{opened} lies on the route by which the email works. Conditioning
    on it removes part of the very effect being measured, and the total effect is no longer
    identified. Do not adjust --- and note that ``opened'' is a variable every marketing
    warehouse holds, which is what makes this the most tempting error in the list.
    Chapter~6 shows how to decompose an effect into direct and indirect parts when the
    decomposition is genuinely what is wanted.
  \item[Collider --- a common effect: $T \rightarrow C \leftarrow Y$.]
    A path through a collider is \emph{blocked by default} and is \emph{opened} by conditioning
    on it. Conditioning on a common effect manufactures association where none existed.
\end{description}

\subsection*{The collider, priced}

The most common self-inflicted wound in marketing analytics is analysing only the customers who
responded --- opened, clicked, replied, converted. \Cref{eq:fnd:collider} makes the damage
visible by planting an email whose true effect is \textbf{exactly zero} and randomizing it, so
that even the naive comparison is unbiased. On all \nbZeroColliderN{} customers, the estimated
email coefficient is \eur{\nbZeroColliderAll} --- correctly, nothing. Restricted to the
\pc{\nbZeroColliderShare} who responded, it is \eur{\nbZeroColliderResp}.

An effect out of thin air, and with the wrong sign. The mechanism is worth stating exactly,
because it is not intuitive: among responders, \emph{having been emailed already explains the
response}, so an emailed responder needed \emph{less} spend to cross the response threshold than
an un-emailed one did. Conditioning on the common effect induces a negative association between
its two causes --- what \citet{pearl2009} calls \emph{explaining away}. Scaled across the
\nbZeroColliderNResp{} responders, the responders-only analysis reports
\eur{\nbZeroColliderPhantom} of phantom revenue destruction on a campaign whose true effect is
zero. Nothing in that analysis looks broken: the treatment was randomized, the sample is large,
the coefficient is precisely estimated. It is simply an answer to a question nobody asked.

The rule this yields is short and absolute: \textbf{post-treatment variables stay out of the
adjustment set}. Opens, clicks, conversions and responses are all descendants of the treatment,
and \cref{as:fnd:backdoor} excludes them.

\subsection*{M-bias: the collider you cannot see}

The uncomfortable corollary --- and the reason \cref{as:fnd:backdoor} is only \emph{partially}
testable --- is that a collider need not be a variable one thinks of as post-treatment, and need
not even be visible in the drawn graph. Consider the classic M-shaped structure, with two
unobserved variables $U_1$ and $U_2$ and one observed pre-treatment covariate $Z$:
%
\begin{equation}
  U_1 \rightarrow Z \leftarrow U_2, \qquad U_1 \rightarrow T, \qquad U_2 \rightarrow Y .
  \label{eq:fnd:mbias}
\end{equation}
%
Here $Z$ is measured \emph{before} treatment, is correlated with both $T$ and $Y$, and passes
every screening rule a practitioner is likely to apply. And yet the path
$T \leftarrow U_1 \rightarrow Z \leftarrow U_2 \rightarrow Y$ is \emph{blocked} --- $Z$ is a
collider on it --- so the naive comparison is unconfounded, and conditioning on $Z$
\emph{opens} the path and \emph{creates} bias. Adjusting for a pre-treatment covariate can make
things worse.

Two things follow. First, ``adjust for everything measured before treatment'' is not a safe
rule; it is merely a safer one than adjusting for everything. Second, the graph is not a picture
of the data --- it is a claim about the world, defended with domain knowledge, and there is no
test on the observed variables that can distinguish \cref{eq:fnd:mbias} from a world in which
$Z$ is an honest confounder. Chapter~7 is the chapter devoted to that problem, and it
demonstrates M-bias with a working example. This chapter's obligation is only to say plainly
that the failure exists, that it is invisible, and that it is why the DAG is drawn first.

% =====================================================================================
\section{The classical baseline: the estimator ladder}
\label{sec:fnd:classical}

The discipline of this book is to ask, before any sampler, what a competent analyst produces
\emph{without} one: a point estimate, and an interval whose covariance assumption is stated out
loud. This section is that read, in full, on the one-confounder world of
\crefrange{eq:fnd:dag_t}{eq:fnd:dag_y}. There are no priors here, no likelihood and no chains.
Chapter~2 adds them --- and its first finding is that they change nothing.

Let $\hat m(t, x) \approx \E[Y \mid T{=}t, X{=}x]$ be a fitted outcome model,
$\hat e(x) \approx \Prob(T{=}1 \mid X{=}x)$ a fitted propensity, and $\bar Y_1, \bar Y_0$ the raw
arm means over $n$ customers. The four rungs are
%
\begin{align}
  \hat\tau_{\text{naive}} &\;=\; \bar Y_1 - \bar Y_0, \label{eq:fnd:naive}\\[2pt]
  \hat\tau_{\text{reg}}   &\;=\; \frac1n \sum_{i=1}^{n}
     \big[\hat m(1, x_i) - \hat m(0, x_i)\big], \label{eq:fnd:reg}\\[2pt]
  \hat\tau_{\text{IPW}}   &\;=\; \frac1n \sum_{i=1}^{n}
     \left[\frac{T_i\,Y_i}{\hat e(x_i)} - \frac{(1 - T_i)\,Y_i}{1 - \hat e(x_i)}\right],
     \label{eq:fnd:ipw}\\[2pt]
  \hat\tau_{\text{AIPW}}  &\;=\; \frac1n \sum_{i=1}^{n}
     \left[\hat m(1, x_i) - \hat m(0, x_i)
      + \frac{T_i\big(Y_i - \hat m(1, x_i)\big)}{\hat e(x_i)}
      - \frac{(1 - T_i)\big(Y_i - \hat m(0, x_i)\big)}{1 - \hat e(x_i)}\right].
     \label{eq:fnd:aipw}
\end{align}

\Cref{eq:fnd:naive} ignores $X$ and is the dashboard. \Cref{eq:fnd:reg} is
\cref{eq:fnd:gformula} with the strata replaced by a fitted line: correct if the outcome model
is correct. \Cref{eq:fnd:ipw} is the same idea from the other side --- weighting a treated
customer by $1/\hat e(x_i)$ makes rarely-emailed kinds of customer count for more, so the
reweighted treated arm matches the \emph{population's} feature mix rather than the treated arm's
\citep{rosenbaum1983}; correct if the propensity model is correct.

\Cref{eq:fnd:aipw} is the doubly-robust estimator of \citet{robins1994}, and its logic deserves
two lines rather than a citation. It is the regression answer \emph{plus} a weighted-residual
correction. If $\hat m$ is right, the residuals $Y - \hat m$ are mean-zero noise, the correction
term vanishes in expectation, and one keeps \cref{eq:fnd:reg}. If instead $\hat e$ is right, the
correction is an unbiased inverse-propensity estimate of precisely the error a wrong $\hat m$
makes, and repairs it. \textbf{Either model being correct suffices} --- hence \emph{doubly}
robust. The implementation here is cross-fitted \citep{chernozhukov2018}: $\hat m$ and $\hat e$
are fitted on one half of the data and applied to the other, so that no customer is scored by a
model that has seen it, which removes an overfitting bias in the correction term.

\subsection*{Each rung, and the standard error it demands}

The standard-error column of \cref{tab:fnd:ladder} is not bookkeeping. An interval built on the
wrong standard error is worse than no interval, because it is a confident lie, and the four rungs
demand four different answers:

\begin{itemize}[leftmargin=*]
  \item \textbf{Naive.} A \emph{Welch} standard error, which does not assume the two arms have
        equal variances --- and here they do not, because loyal customers are both more likely to
        be emailed and noisier spenders.
  \item \textbf{Regression.} An \textbf{HC1} heteroskedasticity-robust standard error
        \citep{white1980}. Spend is noisier at high loyalty; a classical iid standard error would
        quietly understate that.
  \item \textbf{IPW.} No closed form survives the fact that the propensity is itself
        \emph{estimated}, so the interval is a percentile \textbf{bootstrap}
        \citep{efron1979, efron1993}: resample customers with replacement, refit the propensity
        and the estimator, take percentiles of the re-estimates.
  \item \textbf{AIPW.} A cross-fitted influence-function estimate, with a bootstrap interval. It
        is a \emph{frequentist} estimator and its interval is a confidence interval --- a
        distinction Chapter~2 is entirely about.
\end{itemize}

\input{build/tables/nb00_ladder}

\Cref{tab:fnd:ladder} is the chapter's central result. On one confounded sample of
\nbZeroNCustomers{} customers with a planted effect of \eur{\nbZeroTrueDag}, the naive rung
returns \eur{\nbZeroLadNaive} --- \eur{\nbZeroLadNaiveErr}, and its 90\,\% interval,
$[\eur{\nbZeroLadNaiveLo},\ \eur{\nbZeroLadNaiveHi}]$, does not come close to containing the
truth. Regression adjustment returns \eur{\nbZeroLadReg}, inverse-propensity weighting
\eur{\nbZeroLadIpw}, and the doubly-robust estimator \eur{\nbZeroLadAipw}. Every adjusted rung
lands within \eur{\nbZeroLadAdjWorst} of the planted effect. Same data, four estimators, and the
one that ignores the graph is the only one that fails: \textbf{the estimator is the
identification strategy}, and the arithmetic is a detail by comparison.

One detail in \cref{tab:fnd:ladder} is worth pausing on rather than hiding, because it is the
single most misread feature of an interval. On this seed the AIPW interval,
$[\eur{\nbZeroLadAipwLo},\ \eur{\nbZeroLadAipwHi}]$, \emph{just misses} \eur{\nbZeroTrueDag}. That
is not a defect. A 90\,\% interval is \emph{supposed} to miss about one time in ten, and a
procedure whose intervals never missed would be one whose intervals were too wide to be useful.
The correct response to a single miss is not to distrust the estimator but to look at the
distribution of many --- which is the next section, and which is also the exact reason
Chapter~2 spends its length on what an interval does and does not license.

% =====================================================================================
\section{Diagnostics and validation}
\label{sec:fnd:valid}

\subsection*{Bias is what survives repetition}

An estimator is not graded on one draw of the world. \emph{Unbiased} is a statement about the
average over repeated samples, and the honest test refits everything on fresh data.
\Cref{tab:fnd:multiseed} does exactly that, on \nbZeroNSeeds{} fresh confounded panels drawn from
\crefrange{eq:fnd:dag_t}{eq:fnd:dag_y}, and it separates the two failure modes that
\cref{sec:fnd:ident} named:

\begin{itemize}[leftmargin=*]
  \item The naive rung averages \eur{\nbZeroMsNaiveMean} across panels --- a bias of
        \eur{\nbZeroMsNaiveBias} that shows no sign of shrinking, because it is not sampling
        error. It is the confounding, and confounding is a property of the design, not of $n$.
  \item Every adjusted rung centres on the truth: regression at \eur{\nbZeroMsRegMean}, IPW at
        \eur{\nbZeroMsIpwMean}, AIPW at \eur{\nbZeroMsAipwMean}, against a planted
        \eur{\nbZeroTrueDag}. Their sample-to-sample standard deviations
        (\eur{\nbZeroMsRegSd}, \eur{\nbZeroMsIpwSd}, \eur{\nbZeroMsAipwSd}) price what they
        cost in precision, and IPW --- which throws away information by not modelling the
        outcome --- is measurably the noisiest of the three.
  \item The AIPW 90\,\% confidence interval covered the planted effect on \nbZeroAipwCovered{}
        of \nbZeroNSeeds{} panels: \pc{\nbZeroAipwCovPct}, against a nominal 90. The procedure
        keeps its promise. This number is the \emph{only} audit a confidence interval admits, and
        Chapter~2 explains why.
\end{itemize}

\input{build/tables/nb00_multiseed}
\input{build/floats/nb00_rungs}

\subsection*{Overlap and balance: the assumption that can be tested}

\Cref{as:fnd:overlap} is testable, so it is tested. \Cref{fig:fnd:overlap} plots the estimated
propensity by arm. Both arms span essentially the same range, $[\nbZeroEMin, \nbZeroEMax]$, and
the conventional $[0.05, 0.95]$ clip --- which bounds the weights away from division by zero ---
binds on only \pc{\nbZeroClipFrac} of customers. Overlap holds; the IPW rung is trustworthy here.

The companion diagnostic asks whether the weights actually did their job. The
\emph{standardized mean difference} (SMD) in loyalty between the arms --- the gap in means,
measured in pooled standard deviations --- is \nbZeroSmdRaw{} in the raw data, far above the
conventional $0.1$ ``balanced'' bar. Under the inverse-propensity weights it collapses to
\nbZeroSmdIpw. That collapse \emph{is} what an IPW estimator does when it works, and it is the
diagnostic to run before believing one.

\input{build/floats/nb00_overlap}

\begin{remark}
Note the asymmetry the two diagnostics reveal, because it is the shape of the whole field.
Positivity and balance can be checked, and were. Unconfoundedness cannot be checked at all, and
is the assumption everything actually rests on. A causal analysis that reports only the testable
diagnostics has reported only the assumptions it was not worried about.
\end{remark}

% =====================================================================================
\section{What can go wrong}
\label{sec:fnd:wrong}

Everything so far has been a success story, and that is a dangerous thing to teach: the ladder
recovered \eur{\nbZeroTrueDag} because the one confounder was \emph{in the data}. The foundational
question is what happens when it is not.

\subsection*{The unmeasured backdoor}

\Cref{fig:fnd:hidden} and \cref{tab:fnd:hidden} rerun the entire ladder on data in which a hidden
trait $U$ --- a salesperson's hunch, an unlogged loyalty signal --- is wired into both the
assignment and the outcome, while every rung is fitted on the five \emph{observed} features only,
exactly as a real analyst would have to. The naive rung returns \eur{\nbZeroHidNaive} against a
truth of \eur{\nbZeroHiddenTrue}. Regression returns \eur{\nbZeroHidReg}, IPW
\eur{\nbZeroHidIpw}, and AIPW \eur{\nbZeroHidAipw}.

Read that row again. Every adjusted rung sits at least \eur{\nbZeroHiddenMinBias} above the truth,
and \emph{they fail together}. This is the single most important negative result in the chapter,
and it has two parts:

\begin{enumerate}[leftmargin=*]
  \item \textbf{Double robustness offers no protection here}, and the reason is precise rather
        than empirical. AIPW insures against a mis-specified \emph{model of the observed
        features}. It has no insurance to offer against a feature that is \emph{absent}: both of
        its two models are fitted on the same $X$, and the open backdoor runs through $U$, which
        neither can see. Two chances at the wrong problem are not one chance at the right one.
  \item \textbf{The failure is silent.} Nothing in the output of any rung says so. The estimates
        are stable across seeds, the standard errors are small, the diagnostics of
        \cref{sec:fnd:valid} pass. A confidently wrong answer is the default output of a
        well-run analysis on a badly-drawn graph, and Chapter~2 shows that a Bayesian
        posterior fails in exactly the same silent way, and with the same serenity.
\end{enumerate}

\input{build/floats/nb00_hidden_u}
\input{build/tables/nb00_hidden}

The honest responses are three, and none of them is a better estimator: get better data; get a
better \emph{design} (randomize --- \cref{eq:fnd:rct}); or run a \emph{sensitivity analysis},
asking how strong a hidden $U$ would have to be to overturn the decision. Chapter~3 prices
that stress as a tipping point and an E-value.

\subsection*{When positivity fails}

The second assumption fails differently, and worse. Keeping the confounder fully observed but
sharpening the selection until it is nearly deterministic --- so that high-loyalty customers are
almost always emailed and low-loyalty ones almost never --- the propensity clip begins to bind on
\pc{\nbZeroPosClipBad} of customers, and IPW drifts from \eur{\nbZeroPosIpwOk} under gentle
selection to \eur{\nbZeroPosIpwBad} against a truth of \eur{\nbZeroTrueDag}. The estimator is not
being asked to model something badly; it is being asked to compare customers who have no
counterparts in the other arm. No reweighting recovers a full-population ATE from data that
contains no such comparison, and the correct response is to change the estimand --- to report an
effect on the region where overlap holds --- rather than to extrapolate over the region where it
does not. \Cref{fig:fnd:overlap} is the diagnostic that catches this before the number is
trusted, and \cref{sec:fnd:real} is what it looks like when nobody ran it.

% =====================================================================================
\section{On real data: LaLonde}
\label{sec:fnd:real}

Everything above was graded against a truth that was planted. The obvious objection --- that
recovering an effect somebody put there proves nothing --- deserves a real answer, and causal
inference has a canonical one.

\citet{lalonde1986} studied a 1970s job-training programme, the National Supported Work
demonstration. Its decisive feature is that it \emph{included a randomized arm}, so the true
effect of training on 1978 earnings is known: about \usd{\nbZeroLlBenchmark}. LaLonde then asked
the question that founded a literature. Throw the experimental controls away. Replace them with a
\emph{non-experimental} comparison group pulled from national surveys --- the kind of comparison
group an analyst without an experiment actually has --- and see whether the econometric methods of
the day can recover the experimental answer.

They could not, and the failure is spectacular. On the \nbZeroLlTreated{} trainees against
\nbZeroLlControl{} survey controls, the naive difference of means is \usd{\nbZeroLlNaive}. Not
small. \emph{Negative}. Read literally, job training destroyed earnings --- and the reason is
\cref{eq:fnd:decomp} in the wild: the trainees were far more disadvantaged to begin with, so the
selection-bias term swamps the effect and reverses its sign.

Adjusting for the \nbZeroLlNCovs{} measured covariates moves the estimate by \usd{\nbZeroLlSwing},
across zero, to \usd{\nbZeroLlAdj} --- right sign, right order of magnitude, and within
\usd{\nbZeroLlShortfall} of the randomized benchmark (\cref{tab:fnd:lalonde}). This is the whole
chapter, on data nobody simulated: \textbf{the naive comparison is not merely imprecise, it points
the wrong way, and adjustment repairs the sign}.

\input{build/tables/nb00_lalonde}

\subsection*{The residual, and why it is there}

The remaining \usd{\nbZeroLlShortfall} is not a rounding error, and it is not left unexplained.
\Cref{fig:fnd:lalonde} runs \cref{as:fnd:overlap}'s diagnostic --- the same one that passed in
\cref{fig:fnd:overlap} --- on LaLonde, and it fails loudly:
\pc{\nbZeroLlCtrlLowE} of the survey controls sit below $e(x) = 0.10$, against
\pc{\nbZeroLlTrtLowE} of the trainees. Most of the comparison pool looks nothing like a trainee.
A regression adjusting for covariates across that pool is not comparing like with like; it is
\emph{extrapolating} a fitted surface into a region where one arm is essentially absent, and
\cref{eq:fnd:gformula} was never licensed there.

\input{build/floats/nb00_lalonde_overlap}

That is why the field moved to propensity trimming and matching \citep{dehejia1999} --- restrict
the comparison group to the controls that could plausibly have been trainees, and the shortfall
closes --- and why the doubly-robust rung of \cref{eq:fnd:aipw} exists at all. The lesson is not
that adjustment failed on LaLonde. It is that adjustment did most of the job, its own diagnostic
correctly flagged the part it could not do, and the diagnostic pointed at the exact repair.

\begin{remark}[Why this dataset, and not a marketing one]
A marketing analyst may reasonably ask what a job-training programme has to do with an email
campaign. The answer is that LaLonde is the only kind of dataset that can settle the argument: it
carries both a randomized benchmark and an observational comparison group, on the same treatment,
in the same population. That is the exam a causal method has to sit, and real marketing data ---
which has the observational arm and never the benchmark --- cannot set it. Chapter~9 sits the
same exam on a real, randomized geo experiment, and the shape of the finding is identical.
\end{remark}

% =====================================================================================
\section{The decision, in euros}
\label{sec:fnd:euro}

Statistical bias becomes a business problem the moment somebody acts on it. The simulator knows
every customer's true effect $\tau_i$, so each policy can be priced on the actual customer base.

\input{build/tables/nb00_policy}

The confounded read-out --- \eur{\nbZeroNaiveObs} per contact against an \eur{\nbZeroCost} cost ---
says \emph{email everyone}. Doing so on this base realises \eur{\nbZeroPolEveryone}: it burns
\eur{\nbZeroPolBurn}. The true ATE of \eur{\nbZeroAte} sits below the cost and says \emph{email no
one}, which realises nothing and loses nothing, and is the right answer at the level of an
average. \textbf{The identification error is not a footnote; it is a profit-and-loss line}, and its
size here is roughly \eur{\nbZeroPolBurn} per \nbZeroNCustomers{} contacts.

The third row is the one to take away, and it points at the rest of the book. An oracle who could
see each customer's true $\tau_i$ and email only those it pays for would realise
\eur{\nbZeroPolOracle}. That money is not in the average at all --- no ATE-level rule, however
perfectly identified, can reach it --- and it is exactly the prize Chapter~3 competes for.
The decision this chapter fixes is \emph{whether to send}; the decision it cannot reach is
\emph{whom to send to}.

\begin{remark}[The memo to the CMO]
The dashboard says the email pays for itself; it does not. Most of the gap is \emph{who we
emailed}, not what the email did: our historical targeting picked loyal, engaged customers who
would have spent more anyway (\eur{\nbZeroSelBias} of selection bias), and on whom the email also
happens to work better than average (\eur{\nbZeroTargetingGap} of ATT-above-ATE). What to do, in
order: \textbf{randomize} when you can, because a holdout kills selection bias by design; when you
cannot, \textbf{adjust} for the confounders you measured and \textbf{stress-test} the ones you did
not; and note that the real money is not in the average at all.
\end{remark}

% =====================================================================================
\section{Takeaways}
\label{sec:fnd:take}

\begin{enumerate}[leftmargin=*, itemsep=.45em]
  \item \textbf{Half of every unit's table is missing, and that is the whole problem.} Each
        customer has two potential outcomes and reveals one. Causal inference is imputation of
        the other half under stated assumptions \citep{holland1986, ding2018}, and the assumptions
        --- not the arithmetic --- are where the argument is won or lost.

  \item \textbf{The naive difference is an effect plus a claim about who was chosen.}
        \Cref{eq:fnd:decomp} splits it exactly: on this data the dashboard's
        \eur{\nbZeroNaiveObs} is an ATT of \eur{\nbZeroAtt} plus \eur{\nbZeroSelBias} of selection
        bias, and the ATT itself exceeds the ATE by \eur{\nbZeroTargetingGap} because past
        targeting favoured the customers the email works on. Two distinct errors, and neither is
        noise.

  \item \textbf{Randomization is worth more than any estimator in this book.} A coin flip makes
        the selection term exactly zero. The randomized naive estimate is \nbZeroRandSeDist{}
        standard errors from the truth (noise, which shrinks); the observational one is
        \nbZeroObsSeDist{} (bias, which does not). More data makes a confounded estimate more
        precisely wrong.

  \item \textbf{The graph decides what to control for, and two of the three configurations say
        ``don't''.} Confounders are adjusted for; mediators and colliders are not. Conditioning on
        response --- a common effect of the email and of spend --- manufactures
        \eur{\nbZeroColliderResp} per customer of email ``effect'' on a campaign whose true effect
        is \emph{exactly zero}, or \eur{\nbZeroColliderPhantom} of phantom revenue across the
        responders. Post-treatment variables stay out of the adjustment set, and M-bias
        (\cref{eq:fnd:mbias}) shows that even a pre-treatment variable can be a trap.

  \item \textbf{Every adjusted rung works, and the standard error is part of the estimator.} On
        confounded data with a planted \eur{\nbZeroTrueDag}, regression gives \eur{\nbZeroLadReg},
        IPW \eur{\nbZeroLadIpw} and AIPW \eur{\nbZeroLadAipw}, against the naive rung's
        \eur{\nbZeroLadNaive}. Each buys its correctness with a different assumption
        (\cref{tab:fnd:ladder}) and each demands a different covariance --- Welch, HC1, bootstrap.
        An interval built on the wrong standard error is a confident lie.

  \item \textbf{No estimator upgrade fixes missing data.} With an unmeasured confounder, all four
        rungs fail \emph{together} --- every adjusted one at least \eur{\nbZeroHiddenMinBias} above
        the truth --- and they fail \emph{silently}. Double robustness insures against a wrong
        model of the features you have, never against a feature you do not have. Bayes, ML and a
        bigger sample all leave this untouched: \textbf{estimation is sharpened by better methods;
        identification is not}.

  \item \textbf{The failure is real, and it is expensive.} Acting on the confounded
        \eur{\nbZeroNaiveObs} burns \eur{\nbZeroPolBurn} per \nbZeroNCustomers{} contacts. On
        LaLonde it is worse than expensive, it is inverted: the naive comparison says training
        cost workers \usd{\nbZeroLlNaive}, and adjustment moves it \usd{\nbZeroLlSwing} across zero
        to within \usd{\nbZeroLlShortfall} of the experimental truth --- while the overlap
        diagnostic correctly flags the residual it could not close.
\end{enumerate}

\notebooknote{%
  This chapter is \code{notebooks/00\_foundations.ipynb}, \S1--\S6 and \S8--\S9; Chapter~2 is
  its \S7. The simulators are \code{cmp.dgp.uplift\_customers} (\cref{eq:fnd:uplift}) and
  \code{cmp.dgp.dag\_control\_demo} (\crefrange{eq:fnd:dag_t}{eq:fnd:dag_y}); the classical
  estimators and their covariance choices are \code{cmp.classical} and \code{cmp.estimators};
  the seeds are fixed in the notebook. Every number in this chapter comes from the full run
  (\code{CMP\_FAST=0}), with the LaLonde section enabled (\code{CMP\_REAL=1}, which fetches the
  Dehejia--Wahba extract of the NSW data from a public URL), emitted by \code{cmp.report} and
  injected here as a macro. The one exception is the \usd{\nbZeroLlBenchmark} experimental
  benchmark, which is cited from \citet{lalonde1986} rather than estimated, and is flagged as such
  where it is emitted.}
